\chapter{Limitazioni e sviluppi futuri}
\label{ch:limitazioni}

Una presentazione onesta dei risultati richiede di dichiararne i limiti con
la stessa chiarezza con cui se ne espongono i pregi. Questo capitolo
raccoglie le limitazioni del lavoro, distinguendo quelle statistiche, di
esecuzione e strutturali, e indica le direzioni di ricerca che ne
discendono.

\section{Limitazioni statistiche}
\label{sec:lim-statistiche}

La limitazione più rilevante è la scarsa numerosità del campione Bitcoin: 35
trade in-sample e 10 out-of-sample. Gli intervalli di confidenza, per quanto
non includano lo zero, sono ampi, e l'analisi per regimi out-of-sample, con
così pochi trade, non è interpretabile. La validazione cross-asset su SPY
(Capitolo~\ref{ch:validazione}) è stata concepita proprio per mitigare
questo limite, ma non lo elimina sul mercato di interesse primario.

La seconda limitazione è che il periodo out-of-sample è breve e favorevole:
un mercato rialzista con volatilità in compressione gonfia verosimilmente lo
Sharpe out-of-sample rispetto a quello che ci si attenderebbe in un regime
neutro. Per questo, nella tesi, si è dato maggior peso al risultato
in-sample, che comprende invece l'episodio di stress del 2022.

La terza limitazione è l'assenza di un test formale di \emph{data snooping}
sulla strategia finale. Sebbene il percorso di ricerca
(Capitolo~\ref{ch:percorso}) documenti in modo trasparente le ipotesi
scartate, non è stato applicato un test di superiorità predittiva del tipo
di quello di \citet{hansen2005test}, che correggerebbe formalmente per la
molteplicità delle strategie considerate.

\section{Limitazioni di esecuzione}
\label{sec:lim-esecuzione}

Sul piano dell'esecuzione, la copertura delta è ribilanciata con cadenza
giornaliera, mentre Bitcoin presenta rischio di gamma e di salto anche
infragiornaliero: la copertura giornaliera lascia dunque un'esposizione
residua tra un ribilanciamento e il successivo. In secondo luogo, i prezzi
di esecuzione, pur modellati in modo conservativo (ingresso al bid, uscita
all'ask), restano idealizzati: il modello non cattura i riempimenti parziali
né lo slippage in condizioni di mercato sottili. In terzo luogo, la
protezione dal rischio di coda non è implementata, e il suo costo non è
incluso: per una strategia short volatility, l'acquisto di protezione contro
i crash è un costo reale, che una versione operativa dovrebbe internalizzare.

\section{Limitazioni strutturali}
\label{sec:lim-strutturali}

Sul piano strutturale, il vincolo di posizione singola limita la capacità
della strategia: una sola posizione alla volta su un singolo sottostante.
Come mostrato nel Capitolo~\ref{ch:percorso}, l'allentamento di questo
vincolo è stato testato e respinto per via del drawdown inaccettabile, ma
resta il fatto che la strategia ha una capacità di impiego limitata. Inoltre,
il backtest non verifica in ogni istante la sufficienza del margine Deribit:
i margini sono stati calcolati a posteriori (Capitolo~\ref{ch:money-management})
ma non integrati come vincolo dinamico nella simulazione. Infine, l'assenza
di stop loss, giustificata empiricamente, implica che la strategia richiede
la tolleranza a drawdown intraposizione significativi: non è adatta a un
investitore che non possa sostenerli.

\section{Sviluppi futuri}
\label{sec:lim-sviluppi}

Le limitazioni indicano altrettante direzioni di ricerca. Sul fronte del
segnale, il fallimento del gap RND--KDE non chiude la questione del
\emph{timing} del premio: un modello GARCH condizionale sulla misura fisica,
o un modello HAR-RV \citep{corsi2009simple} che sfrutti la memoria lunga
della volatilità realizzata, potrebbero riabilitare un segnale di timing là
dove il confronto pointwise tra densità ha fallito; anche un ensemble di
previsioni di volatilità, che in letteratura tende a battere i singoli
modelli, merita di essere esplorato.

Sul fronte del rilevamento dei regimi, con un campione più ampio
diventerebbero testabili modelli a commutazione di regime di tipo Markov,
l'uso della volatilità della volatilità come indicatore di instabilità, e
l'impiego del premio stesso come predittore dei propri rendimenti futuri,
sulla scia di \citet{bollerslev2009expected}.

Sul fronte dell'esecuzione, una versione operativa dovrebbe incorporare una
protezione dal rischio di coda calcolata dalla superficie di volatilità, una
copertura più frequente e un controllo dinamico del margine.

Infine, sul fronte delle estensioni, l'applicazione a Ethereum, il secondo
mercato di opzioni per liquidità su Deribit, è naturale una volta disponibili
dati storici adeguati. La decomposizione cross-asset del
Capitolo~\ref{ch:validazione} suggerisce inoltre una direzione controintuitiva:
poiché su SPY la gamba opzionaria è profittevole ma la copertura è un costo
netto, una struttura auto-coperta (per esempio uno strangle senza copertura
delta) potrebbe rendere la strategia profittevole anche su asset a bassa
volatilità, eliminando la componente che ne erode il rendimento.
